% Pagina originale 1
\pagina{1}
\tit{Legenda}
\begin{itemize}
\item 1. Insiemi e numeri complessi
\item Minorante e maggiorante
\item Insieme limitato
\item Insieme illimitato
\item Massimo e minimo
\item 1 = Unicità del minimo
\item Intervallo
\item Operazioni tra insiemi
\item Estremo superiore e inferiore
\item 2 = inf (a, b) = a $\land$ sup (a, b) = b
\item Densità di Q in R e di R $\setminus$ Q in R
\item El. Neutro e simmetrico in C
\item Rappresentazione cartesiana di un numero complesso
\item Modulo di un numero complesso
\item Complesso coniugato
\item Rappresentazione polare
\item Argomento e argomento principale
\item Prodotto in forma polare
\item Potenza (Formula di De Moivre)
\item Forma esponenziale (Formula di Eulero)
\item Radice n-esima di un numero complesso
\item Oss. sulla forma esponenziale
\item Oss. sul modulo
\item Disuguaglianza triangolare
\item 2. Funzioni
\item Immagine di una funzione
\item Controimmagine di una funzione
\item Suriettività
\item Iniettività
\item Biettività o invertibilità
\item Funzione inversa
\item Restrizione dell’insieme di partenza
\item Funzione composta
\item Invertibilità di una funzione
\item 3 = $f$ invertibile $\leftrightarrow f$ bigettiva
\item Operazioni tra funzioni
\item Tipi di funzione
\item Funzione potenza
\item Monomio di una variabile
\item Polinomio
\item Funzione lineare e affine
\item Zeri di una funzione
\item Reciproco di una funzione
\item Rapporto tra funzioni
\item Funzione razionale
\item Funzione pari
\item Funzione dispari
\item Funzione periodica
\item Monotonia
\item 4 = $f$ strettamente monotona $\Rightarrow f$ iniettiva
\item 5 = $f,g$ monotone $\Rightarrow f+g$ monotona
\item $f,g$ monotone con stesso tipo $\Rightarrow fg$ monotona
\item 6 = $f$ monotona $\Rightarrow 1/f$ monotona con tipo opposto
\item 7 = $f$ strettamente monotona $\Rightarrow f^{-1}$ strettamente monotona
\end{itemize}
% Pagina originale 2
\pagina{2}
\tit{Legenda}
\begin{itemize}
\item 8 = Monotonia f. composte
\item Successioni
\item Successione estratta o sottosuccessione
\item Monotonia nelle successioni
\item Maggiorante di una funzione
\item Minorante di una funzione
\item Funzione limitata
\item Relazione d’ordine tra funzioni
\item Modulo di una funzione
\item Massimo di una funzione
\item Massimo forte
\item Massimo relativo o locale
\item Inf e sup di una funzione
\item Intorno
\item P.to di accumulazione
\item P.to isolato
\item $+\infty$ punto di accumulazione in $I=(a,+\infty)$ con $a\in\mathbb R$
\item 3. Limiti
\item Definizione algebrica
\item Continuità
\item Oss. sulla continuità
\item f continua in un insieme
\item Limiti nelle successioni
\item 9 = $\{x_n\}$ monotona $\Rightarrow\exists\lim_n x_n=\sup(x_n)\in\overline{\mathbb R}$
\item Convergenza e divergenza nelle successioni
\item Convergenza e divergenza nelle funzioni
\item Infiniti e infinitesimi
\item Teorema ponte
\item $\nexists\lim_{x\to+\infty}\sin x$ usando il teorema ponte
\item Limite da dx e sx
\item P.to di acc. a dx e sx
\item Oss. sui limiti da dx e sx
\item 10 = $\exists\lim_{x\to x_0}f(x)=l\iff\exists\lim_{x\to x_0^-}f(x)=l\land\exists\lim_{x\to x_0^+}f(x)=l$
\item →             →                 →
\item 11 = $\lim_{x\to x_0^+}f(x)=\sup f(x)$
\item Corollario
\item 12 = Th. della permanenza del segno (con limite)
\item 13 = Th. della permanenza del segno (con f. continua)
\item 14 = Conseguenza del th. di permanenza del segno
\item Operazione tra limiti
\item 15 = Teorema degli zeri
\item Oss. th. degli zeri
\item 16 = Conseguenze th. degli zeri
\item Oss. sul corollario
\item 17 = Th. dei valori intermedi
\item 18 = Th. di Weierstrass
\item Th. valori intermedi + Th. di Weierstrass
\item Funzioni elementari
\item Funzioni trigonometriche (e inverse)
\item Funzioni iperboliche (e inverse)
\item 19 = Th. di unicità del limite
\item Rapporto tra limiti
\item $\lim_{x\to x_0}f(x)=l^-,\quad\lim_{x\to x_0}f(x)=l^+$
\item →               →
\item 20 = Th. del doppio confronto
\item 21 = Th. del confronto
\item 22 = $\exists\lim_{x\to x_0}f(x)=0\iff\exists\lim_{x\to x_0}|f(x)|=0$
\item →             →
\end{itemize}
% Pagina originale 3
\pagina{3}
\tit{Legenda}
\begin{itemize}
\item Prodotto tra 1/$\infty$ e limitato
\item 23 = Th. limite di funzioni composte
\item 24 = Th. continuità di f. composte da f. continue
\item Punti di discontinuità
\item Eliminazione p.ti di discontinuità di 3° specie
\item Asintoti
\item Oss. sugli asintoti nelle funzioni simmetriche
\item Simboli di Landau
\item Funzioni asintotiche
\item Limiti notevoli
\item 4. Derivate
\item Funzione dotata di derivata in $x_0$
\item $f_-'(x)\ne f_+'(x)\Rightarrow\nexists f'(x)$
\item $f_-'(x)=f_+'(x)\Rightarrow\exists f'(x)$
\item Miglior approssimazione lineare in un punto
\item 25 = $f$ derivabile in $x_0\iff f$ ha miglior approssimazione lineare in $x_0$
\item Oss. sulla miglior approssimazione lineare
\item 26 = $f$ derivabile $\Rightarrow f$ continua
\item 27 = $f$ derivabile a destra e sinistra $\Rightarrow f$ continua
\item P.ti non derivabili
\item Regole di derivazione
\item 28 = Derivata di f. composta
\item 29 = Derivata della f. inversa
\item Punto stazionario
\item 30 = Th. di Fermat
\item 31 = Th. di Rolle
\item 32 = Th. di Cauchy
\item 33 = Th. di Lagrange
\item 34 = Corollario
\item 35 = Caratterizzazione della monotonia mediante il segno della derivata
\item 36 = $f'(x)>0\Rightarrow f$ strettamente crescente
\item 5. Integrali indefiniti
\item Primitiva di una funzione
\item 37 = Le primitive di f distano di una costante
\item 38 = Linearità dell’integrale indefinito
\item 39 = Integrazione per parti
\item 40 = Integrazione per sostituzione
\item 6. Derivate seconde + De L’Hopital
\item 41 = Th. di De L’Hopital
\item Applicazioni del th. di De L’Hopital
\item 42 = $\exists\lim_{x\to x_0}f'(x)=l\Rightarrow\exists f'(x)=l$
\item Convessità
\item 43 = Conseguenze della convessità
\item 44 = $\exists f'(x_0)\Rightarrow\exists f''(x_0)\in\mathbb R$
\item 45 = $f$ convessa $\iff f''>0$
\item P.to di flesso
\item 46 = $x_0$ punto di flesso $\Rightarrow f''(x_0)=0$
\item Tipi di p.ti di flesso
\item Derivate elementari
\item 7. Serie di Taylor
\item Formula di Taylor
\item 47 = $P_n(x_0)=f(x_0), P_n'(x_0)=f'(x_0),\ldots,P_n^{(k)}(x_0)=f^{(k)}(x_0)$
\item 48 = $f(x)-P_n(x)=o((x-x_0)^n)$
\item 49 = Caratterizzazione punto estremale tramite $f^{(n)}$
\item Polinomio di McLaurin
\end{itemize}
% Pagina originale 4
\pagina{4}
\tit{Legenda}
\begin{itemize}
\item 8. Integrale di Riemann
\item Somma inferiore e superiore
\item Oss. sugli inf e sup
\item 50 = Relazione tra s ed S
\item Integrabilità secondo Riemann
\item Integrale f. costante
\item Continuità $\Rightarrow$ integrabilità
\item Ampiezza di una suddivisione
\item 51 = $f$ crescente $\Rightarrow f$ integrabile
\item Proprietà integrale di Riemann
\item Media integrale
\item 52 = Th. fondamentale del Calcolo integrale
\item Integrazione per parti
\item Integrazione per sostituzione
\item Oss. sulla simmetria
\item 53 = Continuità f. integrale
\end{itemize}
